% Generated by Listas/make_statements.py from pratica-lecture-06.md, item 21. Do not edit.
\textbf{Source.} Cabral (2017), \emph{Introduction to Industrial Organization}, 2nd ed., Ch. 9 Review and practice exercises, Exercise 9.15 ``Multimarket contact'', p.~244 {[}PDF 256{]}. Reproduced verbatim. Answer in the source: none in the book.



\begin{quote}

\textbf{9.15. Multimarket contact.} Consider the model of multimarket contact presented in Section 9.3: Firm 1 has cost \(\underline{c}\) in Market 1, while Firm 2 has a cost \(\overline{c}\). The situation is reversed in Market 2. Demand is the same in both markets: consumers are willing to buy \(q\) units for a price of up to \(p^M\) (that is, for \(p \le p^M\), quantity demanded does not depend on price). It is assumed that \(\underline{c} < \overline{c} < p^M\). In each period, firms set prices in both markets simultaneously.



Determine the minimum value of the discount factor such that the optimal collusive solution is stable.

\end{quote}

